\par\Needspace{6\baselineskip}\medskip
\SourceLeaf{195}\begin{center}\hypertarget{native-195-article144}{}\pdfbookmark[2]{144.}{article-144}144.\end{center}
Arcubus $\zeta,\zeta''$ inventis, punctorum $C,C''$ positio data erit, poteritque distantia
$CC''=2f'$ ex $\zeta,\zeta''$ et $\epsilon'$ determinari. Sint $u,u''$ inclinationes circulorum
maximorum $AB,A''B''$ ad circulum maximum $CC''$ (quae in figura 4 resp.
erunt anguli $C''CD'$ et $180^{0}-CC''D'$), habebimusque aequationes sequentes,
aequationibus 3—6 art. 137 prorsus analogas:
\[
\begin{aligned}
\sin f'\sin\tfrac12(u''+u)&=\sin\tfrac12\epsilon'\sin\tfrac12(\zeta+\zeta'')\\
\sin f'\cos\tfrac12(u''+u)&=\cos\tfrac12\epsilon'\sin\tfrac12(\zeta-\zeta'')\\
\cos f'\sin\tfrac12(u''-u)&=\sin\tfrac12\epsilon'\cos\tfrac12(\zeta+\zeta'')\\
\cos f'\cos\tfrac12(u''-u)&=\cos\tfrac12\epsilon'\cos\tfrac12(\zeta-\zeta'').
\end{aligned}
\]
Duae priores dabunt $\tfrac12(u''+u)$ et $\sin f'$, duae posteriores $\tfrac12(u''-u)$ et $\cos f'$; ex
$\sin f'$ et $\cos f'$ habebitur $f'$. Angulos $\tfrac12(u''+u)$ et $\tfrac12(u''-u)$, qui in ultima demum
hypothesi ad determinandum situm plani orbitae adhibebuntur, in hypothesibus
primis negligere licebit.

Prorsus simili modo $f$ ex $\epsilon,C'D$ et $C''D$, nec non $f''$ ex $\epsilon'',CD'',C'D''$
derivari possent: sed multo commodius ad hunc finem formulae sequentes adhibentur
\[
\begin{aligned}
\sin2f&=r\sin2f'\cdot\frac{n}{n'r'}\\
\sin2f''&=r''\sin2f'\cdot\frac{n''}{n'r'}\mathpunct{\raisebox{.45ex}{,}}
\end{aligned}
\]
ubi logarithmi quantitatum $\dfrac{n}{n'r'}\mathpunct{\raisebox{.45ex}{,}}$ $\dfrac{n''}{n'r'}$ iam e calculis praecedentibus adsunt.
Totus denique calculus confirmationem novam inde nanciscetur, quod fieri
debet $2f+2f''=2f'$: si qua forte differentia prodeat, nullius certe momenti
esse poterit, siquidem omnes operationes quam accuratissime peractae fuerint.
Interdum tamen, calculo ubique septem figuris decimalibus subducto, ad aliquot
minuti secundi partes decimas assurgere poterit, quam si operae pretium
videtur facillimo negotio inter $2f,2f''$ ita dispertiemur, ut logarithmi sinuum
aequaliter vel augeantur vel diminuantur, quo pacto aequationi $P=\dfrac{r\sin2f''}{r''\sin2f}=\dfrac{n''}{n}$
omni quam tabulae permittunt praecisione satisfactum erit. Quoties $f$ et
$f''$ parum differunt, differentiam illam inter $2f$ et $2f''$ aequaliter distribuisse
sufficiet.
